% Resume bullet points -- Warehouse Robotics Simulation
% MathWorks Excellence in Innovation, Project 212
% Framing: autonomous driving / control systems.
% Drop the \item lines into your existing resume's itemize/cvitems block.

%========================================================================
% VERSION A -- Concise (4 bullets), for a tight one-page resume
%========================================================================
\begin{itemize}
  \item Designed a closed-loop autonomous navigation stack in MATLAB for a
        differential-drive mobile robot, integrating motion planning,
        trajectory tracking, and collision avoidance over a simulated
        50\,m\,$\times$\,30\,m environment.
  \item Implemented a pure-pursuit path-tracking controller with tunable
        lookahead and saturated linear/angular velocity commands, driving a
        unicycle vehicle model along planned trajectories with closed-loop
        pose feedback.
  \item Developed and benchmarked decision-making and motion-planning
        algorithms (A*, space-time prioritized planning, and optimal
        Conflict-Based Search) for multi-vehicle coordination, evaluating
        tracking cost, makespan, and runtime as fleet size scaled from 2 to 8.
  \item Validated the control stack in a 3D physics simulator (Gazebo
        co-simulation) by commanding the vehicle through the MATLAB
        co-simulation interface, including differential-drive inverse
        kinematics mapping body velocities to wheel commands.
\end{itemize}

%========================================================================
% VERSION B -- Detailed (6 bullets), for a projects section / CV
%========================================================================
\begin{itemize}
  \item Built a full autonomous-navigation pipeline -- perception map,
        global planner, and motion controller -- in MATLAB for a
        differential-drive robot (MathWorks Excellence in Innovation
        Project~212), modeling the environment as an obstacle-inflated
        occupancy grid for guaranteed clearance.
  \item Implemented a pure-pursuit lateral/longitudinal controller that tracks
        planned waypoints using a tunable lookahead distance and bounded
        linear and angular velocity, closing the loop on the vehicle's pose at
        each timestep over a kinematic unicycle model.
  \item Derived and applied differential-drive inverse kinematics to convert
        commanded body velocity ($v$, $\omega$) into left/right wheel
        angular-velocity setpoints, the actuation layer between the
        controller and the vehicle.
  \item Engineered global path planning with A* and extended it to a
        space-time ($x$,$y$,$t$) planner with a shared reservation table for
        multi-vehicle coordination, resolving vehicle--vehicle vertex and
        head-on swap conflicts.
  \item Implemented Conflict-Based Search (CBS), a complete and optimal
        multi-agent motion planner, and built a benchmark quantifying the
        optimality--compute trade-off versus the prioritized baseline
        (sum-of-costs, makespan, success rate, planning time vs.\ fleet size).
  \item Performed simulation-based validation of the autonomy stack in Gazebo
        via co-simulation, commanding the vehicle through the Robotics System
        Toolbox \texttt{gz} interface; stood up a reproducible headless
        Gazebo~11 environment in Docker on Apple Silicon (arm64 plugin build,
        software-OpenGL rendering over VNC).
\end{itemize}

%========================================================================
% Skills line (optional)
%========================================================================
% \textbf{Skills:} MATLAB, control systems (pure-pursuit path tracking,
% closed-loop feedback), motion planning (A*, multi-agent path finding/CBS),
% differential-drive kinematics, autonomous navigation, Gazebo co-simulation,
% Docker.
